% =============================================================================
%  CHAPTER 4 - RESULTS OF IMPACT TEST ANALYSIS
% =============================================================================
\chapter{Results of Impact Test Analysis}
\label{chap:impact-results}

This chapter reports the bulk-specimen response of the nine-configuration
impact-loading programme. The chapter has three purposes. First, it
documents the settlement history of the four tests in which axial
settlement was monitored at five blow milestones, and uses the
shake-down concept to justify the choice of twenty blows as the loading
endpoint. Second, it reports the peak hammer force imparted to the
specimen at the two drop heights used in the programme, confirming that
the apparatus delivered the nominal energy levels reproducibly across the
nine tests. Third, it presents the post-test gradation analysis of the
specimens through the Indian Railway Fouling Index \FIIN{} and examines
the relationship between settlement and fouling.

\section{Axial settlement}
\label{sec:imp-settlement}

The cumulative axial settlement of the ballast specimen was recorded at
five blow milestones ($N = 0, 5, 10, 15, 20$) in four of the nine tests:
T1 (PUBM-FB-300), T3 (UR-FB-300), T5 (RUBM-RB-450), and T9 (PUBM-RB-300).
The four settlement curves recorded are plotted in
Figure~\ref{fig:imp-settlement-cum}, with each curve referenced to its
pre-first-blow datum so that what is shown is the net cumulative settlement.

\begin{figure}[ht]
    \centering
    \includegraphics[width=0.85\textwidth]{image25.png}
    \caption{Cumulative net axial settlement of the ballast specimen as a
             function of hammer-blow count $N$ for the four monitored
             tests. The shaded band at high $N$ marks the shake-down
             (steady-state) regime within which incremental settlement per
             blow falls below approximately \SI{0.5}{\milli\metre}.}
    \label{fig:imp-settlement-cum}
\end{figure}

\begin{figure}[ht]
    \centering
    \includegraphics[width=0.85\textwidth]{image26.png}
    \caption{Incremental settlement per blow within each five-blow
             interval, for the four monitored tests. The fall in the
             per-blow rate from the \numrange{0}{5} interval to the
             \numrange{15}{20} interval defines the transition from active
             densification to the shake-down regime.}
    \label{fig:imp-settlement-inc}
\end{figure}

Figure~\ref{fig:imp-settlement-inc} establishes that in all four monitored
tests the per-blow incremental settlement at the \numrange{15}{20}
interval has fallen to between approximately
\SI{0.05}{\milli\metre\per\text{blow}} and
\SI{0.5}{\milli\metre\per\text{blow}}---between one and two orders of
magnitude below the rate in the first interval. Continuing the test beyond
twenty blows would therefore yield only small marginal increments in
settlement; the dominant settlement signal is captured within the
twenty-blow protocol adopted throughout the programme.

\section{Peak hammer force per test}
\label{sec:imp-peak-force}

The in-line \SI{5}{\tonne} load cell beneath the striker plate recorded
the peak compressive force imparted on every blow of every test. Across
the nine-configuration programme the peak force grouped tightly by drop
height: tests run at the \SI{300}{\milli\metre} drop produced peak forces
around \SI{19.4}{\kilo\newton} per blow (corresponding to a nominal
contact stress of approximately \SI{275}{\kilo\pascal} over the
\SI{300}{\milli\metre} bearing-plate area), and tests run at the
\SI{450}{\milli\metre} drop produced peak forces around
\SI{37.0}{\kilo\newton} per blow (approximately \SI{523}{\kilo\pascal}).
Figure~\ref{fig:imp-peak-force} plots the average peak force for each of
the nine tests, grouped along the horizontal axis by drop height. The
two-tier structure is consistent across pad type: the presence of a UR,
RUBM or PUBM pad does not appreciably change the peak force at the load
cell because the pad is below the ballast-fill column and the load cell
measures the through-going force at the base of the assembly, not at the
ballast surface.

\begin{figure}[ht]
    \centering
    \includegraphics[width=0.85\textwidth]{image27.png}
    \caption{Mean peak hammer force recorded by the in-line \SI{5}{\tonne}
             load cell for each of the nine tests, grouped by drop height.}
    \label{fig:imp-peak-force}
\end{figure}

\section{Indian Railway Fouling Index \texorpdfstring{$\FIIN$}{FI\_IN}}
\label{sec:imp-FIIN}

The post-test gradation of every specimen was determined by sieve analysis
through the standard RDSO sieve series (63, 40, 31.5, 20, 16, 12.5, 10,
4.75, 2.36, 1.18, 0.6, \SI{0.075}{\milli\metre} and pan). From the
resulting sieve distribution the Indian Railway Fouling Index \FIIN{} was
computed as defined in Section~\ref{sec:meth-indices} and after Anbazhagan
et~al.\ (2012):
\begin{equation}
    \FIIN = \PoneM + \PtwoM,
    \label{eq:imp-FIIN}
\end{equation}
where \PoneM{} is the percentage of post-test mass passing the
\SI{0.075}{\milli\metre} sieve (the fines fraction) and \PtwoM{} is the
percentage passing the \SI{20}{\milli\metre} sieve (the fraction below the
RDSO minimum acceptable particle size). The two terms have different
physical interpretations: \PoneM{} captures the fines that drive drainage
degradation and the loss of inter-particle friction, while \PtwoM{}
captures the coarser fraction that has migrated down out of the
RDSO-acceptable envelope and that, once accumulated in the field, requires
either tamping or undercutting to restore. Both contribute to the overall
ballast fouling state.

\begin{figure}[ht]
    \centering
    \includegraphics[width=0.90\textwidth]{image28.png}
    \caption{Indian Railway Fouling Index $\FIIN = \PoneM + \PtwoM$ (full
             bar) and the \PoneM{} fines contribution (orange overlay) for
             each of the nine tests.}
    \label{fig:imp-FIIN}
\end{figure}

\begin{table}[ht]
\centering
\caption{Indian Railway Fouling Index $\FIIN$, the fines fraction
         $\PoneM$, and the \SI{20}{\milli\metre}-passing fraction $\PtwoM$
         for the nine test configurations. Values computed per the
         $\FIIN$ definition of Section~\ref{sec:meth-indices} from the
         post-test sieve data archived in the laboratory spreadsheet.}
\label{tab:imp-FIIN}
\begin{tabular}{@{}cccccc@{}}
    \toprule
    \textbf{Test} & \textbf{Configuration}
    & $\PoneM$ (\si{\percent})
    & $\PtwoM$ (\si{\percent})
    & $\FIIN$ (\si{\percent}) \\
    \midrule
    T1 & PUBM-FB-300 & 0.027 & 0.507 & 0.533\\
    T2 & RUBM-FB-450 & 0.020 & 1.339 & 1.359\\
    T3 & UR-FB-300   & 0.028 & 1.384 & 1.412\\
    T4 & UR-RB-450   & 0.046 & 1.914 & 1.960\\
    T5 & RUBM-RB-450 & 0.049 & 1.307 & 1.356\\
    T6 & PUBM-RB-450 & 0.017 & 0.747 & 0.764\\
    T7 & UR-RB-450   & 0.118 & 3.268 & 3.386\\
    T8 & RUBM-RB-300 & 0.036 & 0.631 & 0.667\\
    T9 & PUBM-RB-300 & 0.019 & 0.953 & 0.972\\
    \bottomrule
\end{tabular}
\end{table}

Four observations from Figure~\ref{fig:imp-FIIN} and
Table~\ref{tab:imp-FIIN} are worth setting out explicitly. First, the
\FIIN{} values across the nine tests span a factor of approximately
\num{6.4}---from \SI{0.533}{\percent} at T1 (PUBM-FB-300) to
\SI{3.386}{\percent} at T7 (UR-RB-450)---which establishes that the
test-configuration variables (pad type, base condition, drop height)
exert a material influence on the bulk fouling response. Second, the
highest-fouling configuration is T7 (UR-RB-450), the combination of no
pad with a rigid base at the higher drop energy. This is mechanically the
most aggressive corner of the experimental matrix---the impulsive energy
is not attenuated by either a pad or a compliant base---and the result is
consistent with that expectation. Third, the lowest-fouling configuration
is T1 (PUBM-FB-300), the combination of a stiff polyurethane pad with a
compliant sand base at the lower drop energy. The compliant sand base
absorbs a substantial fraction of the impulsive energy through frictional
and re-arrangement losses, and the lower drop reduces the per-blow input
energy by one third relative to the \SI{450}{\milli\metre} drop.

Fourth, the fines contribution \PoneM{} is small in absolute terms across
the programme. It contributes between \SI{0.017}{\percent} (T6,
PUBM-RB-450) and \SI{0.118}{\percent} (T7, UR-RB-450) of the post-test
mass---only a small fraction of the total \FIIN, which is dominated by
the \PtwoM{} term. This indicates that the dominant bulk-level damage
mechanism in the present programme is the migration of particles below
the \SI{20}{\milli\metre} RDSO acceptable lower bound (i.e.\
fragmentation that produces \SI{16}{\milli\metre} and smaller daughter
particles), rather than the production of fines through abrasion of
asperities. The grain-scale diagnostic of
Chapter~\ref{chap:morph-results} will revisit this distinction directly.

A further pattern that emerges from Table~\ref{tab:imp-FIIN} is the
relative ranking of the three pads on each base. On the rigid base at
\SI{450}{\milli\metre} drop the PUBM pad (T6, $\FIIN =
\SI{0.764}{\percent}$) outperforms both the RUBM pad (T5,
\SI{1.356}{\percent}) and the no-pad control (T7,
\SI{3.386}{\percent}); on the rigid base at \SI{300}{\milli\metre} drop
the same ranking holds (T9 PUBM \SI{0.972}{\percent} < T8 RUBM
\SI{0.667}{\percent}), with both pads outperforming the comparable
unreinforced configuration. On the flexible base at
\SI{300}{\milli\metre} drop the PUBM pad (T1, \SI{0.533}{\percent})
outperforms the no-pad control (T3, \SI{1.412}{\percent}), again
confirming that adding a polyurethane mat below the ballast reduces bulk
fouling. The pad--base coupling pattern therefore points consistently
toward the PUBM as the most effective fouling-reducer in the matrix
tested.

\section{Settlement--fouling correlation}
\label{sec:imp-settlement-fouling}

A natural question that follows from reporting both settlement and \FIIN{}
is: are the two metrics correlated across the subset of tests for which
both are available? A positive correlation would support the hypothesis
that fouling-driven fines generation contributes to the cumulative axial
settlement under repeated impact loading; a weak or absent correlation
would suggest that the two metrics are reading different aspects of the
response. The correlation is therefore evaluated on the four-test
settlement subset (Figure~\ref{fig:imp-corr}).

\begin{figure}[ht]
    \centering
    \includegraphics[width=0.80\textwidth]{image29.png}
    \caption{Indian Railway Fouling Index \FIIN{} plotted against the net
             cumulative settlement at $N = 20$ for the four monitored
             tests.}
    \label{fig:imp-corr}
\end{figure}

Two observations follow from Figure~\ref{fig:imp-corr}. First, within the
three low-energy \SI{300}{\milli\metre} tests (T1, T3, T9) a
weakly-positive trend exists between settlement and fouling: the test
with the largest settlement (T3, UR-FB-300 at \SI{38.0}{\milli\metre}) is
also the one with the largest \FIIN{} (\SI{1.412}{\percent}), and the
test with the smallest settlement (T9, PUBM-RB-300 at
\SI{24.6}{\milli\metre}) has an intermediate \FIIN{}
(\SI{0.972}{\percent}). The slope is shallow and the scatter is wide
given only three points, so this trend should be read as directional
rather than predictive. Second, the high-energy test T5 (RUBM-RB-450)
lies clearly off the \SI{300}{\milli\metre} trend: it has the largest
settlement of the four (\SI{48.6}{\milli\metre}) but a comparatively
modest \FIIN{} (\SI{1.356}{\percent}, actually similar to the value at
the much-lower-settlement T2 RUBM-FB-450 in Table~\ref{tab:imp-FIIN}).
Settlement and fouling are therefore not interchangeable bulk metrics;
they index different aspects of the underlying degradation.

The decoupling is mechanistically intelligible. Settlement reflects axial
deformation from particle rearrangement and void closure, which need not
involve fragmentation. Fouling, by contrast, requires actual particle
breakage and is governed by peak contact force and parent-rock
brittleness. Thus T5 settles substantially without much fragmentation,
while T7 fragments without the largest settlement.

\section{Discussion}
\label{sec:imp-discussion}

The bulk findings of
Sections~\ref{sec:imp-settlement}--\ref{sec:imp-settlement-fouling} are
broadly consistent with the published drop-weight impact literature on
ballast. The shock-mat-on-rigid-base condition explored by Nimbalkar,
Indraratna, Dash and Christie (2012) produced a similar pattern of
reduced bulk degradation when a compliant pad was placed beneath the
ballast on a rigid support, and the natural-rubber and recycled-rubber
mats in that study delivered the largest improvement on the stiffer
support---directionally consistent with the present finding that PUBM
substantially outperforms the unreinforced configuration on the rigid
base. The crumb-rubber-modified ballast studies of Koohmishi and
Azarhoosh (2021, 2022) and Zhang et~al.\ (2023) reported FI
fines-generation values in the same order of magnitude as the present
\PoneM{} results at comparable energy, and Koohmishi and Palassi (2020)
reported similar fragmentation severity on basalt ballast under matched
drop-weight loading. The shake-down framework used to justify the
twenty-blow protocol in Section~\ref{sec:imp-settlement} follows the
treatment of Werkmeister et~al.\ (2001) for granular materials and is the
same framework used in the impact-loading context by Nimbalkar et~al.\
(2012).

Three patterns visible at the bulk level are worth recording before
moving to Chapter~\ref{chap:morph-results}. First, the PUBM-on-rigid-base
configurations T6 and T9 outperform the RUBM and UR equivalents at both
drop heights, suggesting that a stiffer polyurethane mat is more
effective than a softer rubber mat at attenuating impulsive loading on a
rigid support. Second, the highest-\FIIN{} configuration is T7
(UR-RB-450), the no-pad rigid-base high-energy combination---the
mechanically most aggressive corner of the experimental matrix and the
one most representative of the unprotected rigid-substructure conditions
encountered at bridge approaches and welded-rail transition zones in
service. Third, the settlement and fouling responses are not
interchangeable: the test with the largest settlement (T5) is not the
test with the largest fouling (T7). Both metrics are needed to
characterise the bulk response, and even together they leave open the
question of which grain-scale processes drive each metric in each
configuration. That is the question addressed in
Chapter~\ref{chap:morph-results}.
